% !TEX program = lualatex
\documentclass[lualatex,ja=standard]{bxjsarticle}

\usepackage{amsmath}
\usepackage{amssymb}   % \mathbb（ℝ, ℕ など）
\usepackage{xcolor}

% -------------------------------------------------------------------
% \newcommand: 引数なし
% -------------------------------------------------------------------
\newcommand{\myname}{LaTeX サンプルユーザー}
\newcommand{\R}{\mathbb{R}}          % 実数集合 ℝ
\newcommand{\N}{\mathbb{N}}          % 自然数集合 ℕ
\newcommand{\highlight}[1]{\textbf{\textcolor{red}{#1}}}  % 1引数

% -------------------------------------------------------------------
% \newcommand: 引数あり
% -------------------------------------------------------------------
% #1 = 分子, #2 = 分母
\newcommand{\myfrac}[2]{\dfrac{#1}{#2}}

% デフォルト引数付き（オプション引数）
% #1 = 省略可能なオプション（デフォルト "blue"）, #2 = テキスト
\newcommand{\colortext}[2][blue]{\textcolor{#1}{#2}}

% -------------------------------------------------------------------
% \newenvironment: 独自環境
% -------------------------------------------------------------------
\newenvironment{mybox}
  {\begin{center}
   \begin{tabular}{|c|}
   \hline
   \begin{minipage}{0.85\linewidth}\smallskip\par\noindent\ignorespaces}
  {\par\smallskip\end{minipage}\\
   \hline\end{tabular}\end{center}}

% -------------------------------------------------------------------
% \renewcommand: 既存コマンドの上書き
% -------------------------------------------------------------------
% \epsilon をより見やすい \varepsilon に変更
\renewcommand{\epsilon}{\varepsilon}

\title{カスタムコマンドのサンプル}
\author{\myname}
\date{\today}

\begin{document}

\maketitle

% -------------------------------------------------------------------
% 引数なしコマンド
% -------------------------------------------------------------------
\section{引数なしコマンド}

著者は \texttt{\textbackslash myname} → \myname

実数集合 $\R$、自然数集合 $\N$（\texttt{\textbackslash R}, \texttt{\textbackslash N}）

% -------------------------------------------------------------------
% 1引数コマンド
% -------------------------------------------------------------------
\section{1引数コマンド}

\texttt{\textbackslash highlight\{テキスト\}} → \highlight{重要な箇所}

% -------------------------------------------------------------------
% 2引数コマンド
% -------------------------------------------------------------------
\section{2引数コマンド}

\texttt{\textbackslash myfrac\{分子\}\{分母\}} → $\myfrac{a+b}{c-d}$

% -------------------------------------------------------------------
% デフォルト引数付きコマンド
% -------------------------------------------------------------------
\section{デフォルト引数付きコマンド}

オプション省略（デフォルト blue）: \colortext{青いテキスト}

オプション指定: \colortext[red]{赤いテキスト}

オプション指定: \colortext[purple]{紫のテキスト}

% -------------------------------------------------------------------
% 独自環境
% -------------------------------------------------------------------
\section{独自環境}

\begin{mybox}
  これは \texttt{mybox} 環境で囲まれたテキストです。
  枠線付きの中央揃えボックスとして表示されます。
\end{mybox}

% -------------------------------------------------------------------
% \renewcommand の効果
% -------------------------------------------------------------------
\section{renewcommand}

\texttt{\textbackslash epsilon} を \texttt{\textbackslash varepsilon} に上書きしたため、
$\epsilon$ が自動的に \(\varepsilon\) の形状になります。

% -------------------------------------------------------------------
% \newcommand vs \providecommand
% -------------------------------------------------------------------
\section{providecommand}

\texttt{\textbackslash providecommand} は、コマンドが未定義のときだけ定義します。
既に定義済みの場合は何もしません。パッケージとの競合回避に便利です。

\providecommand{\etal}{\textit{et al.}}

Smith \etal (2023) はその例です。

\end{document}